\documentclass[10pt,a4paper]{amsart}
\usepackage[margin=19mm]{geometry}
\usepackage{amsmath,amssymb,mathtools}
\usepackage[scr=rsfs,cal=euler]{mathalfa}
\usepackage{xfrac}
\usepackage[unicode,hidelinks,bookmarks=false]{hyperref}
\newcommand{\ie}{i.e.\ }
\newcommand{\cf}{cf.\ }
\newcommand{\resp}{respectively\ }
\newcommand{\Subsection}{\subsection}
\providecommand{\nobreakdash}{\nobreak\hbox{-}}
\setlength{\emergencystretch}{2em}
\multlinegap0pt
\usepackage{mathtools}
\usepackage{booktabs}
\allowdisplaybreaks
\numberwithin{equation}{section}
\newtheorem{theorem}{Theorem}[section]
\newtheorem{lemma}[theorem]{Lemma}
\newtheorem{proposition}[theorem]{Proposition}
\newtheorem{corollary}[theorem]{Corollary}
\newtheorem{conjecture}[theorem]{Conjecture}
\newtheorem{definition}[theorem]{Definition}
\newtheorem{remark}[theorem]{Remark}
\newcommand{\N}{\mathbb{N}}
\newcommand{\Z}{\mathbb{Z}}
\newcommand{\F}{\mathbb{F}}
\newcommand{\eps}{\varepsilon}
\newcommand{\bw}{\operatorname{bw}}
\newcommand{\ceil}[1]{\left\lceil #1\right\rceil}
\newcommand{\floor}[1]{\left\lfloor #1\right\rfloor}
\begin{document}
\title[Arithmetic Progression-Free Subset-Sum Sets]{Arithmetic Progression-Free Subset-Sum Sets}
\author{Samuel Korsky}
\date{}
\maketitle
\noindent\emph{Source and licence.} Arithmetic Progression-Free Subset-Sum Sets. Version arXiv:2606.24139v1 (CC BY 4.0). This material is distributed under the Creative Commons Attribution 4.0 International licence, http://creativecommons.org/licenses/by/4.0/, as recorded in the arXiv OAI licence metadata for this exact version and shown on the abstract page https://arxiv.org/abs/2606.24139v1. Full rights notice: licenses/arxiv-2606.24139-CC-BY-4.0.txt.\par\smallskip
\noindent\textit{Editorial scope.} Counted unit: the original \texttt{theorem} environment \texttt{thm:graph-construction} at source lines 675--690 together with its complete proof at lines 692--709; the delivered document keeps source lines 639--710 so that every referenced label and every citation resolves inside it. Native editable source is the paper's own concatenated TeX; the AMS wrapper is editorial formatting only. No publisher class, upstream script or unrelated artwork is redistributed.
\section{Distinct-Generator Digit Constructions}\label{sec:upper}

We begin with the progression-free digit language.

\begin{lemma}[Restricted digits]\label{lem:restricted-digits}
Let $p$ be prime, let $k\ge3$, and let $0\le\tau\le\min\{p,k\}-2$. The set of nonnegative integers whose base-$p$ digits all belong to $\{0,1,\ldots,\tau\}$ contains no nonconstant $k$-term arithmetic progression.
\end{lemma}

\begin{proof}
Suppose $x_j=x_0+jd$ for $0\le j<k$ and $d>0$. Let $s=v_p(d)$ and let $u=d/p^s\not\equiv0\pmod p$. Looking at the $s$th base-$p$ digit modulo $p$, the digits of $x_j$ are congruent to $c+ju\pmod p$ for some $c$. If $k\le p$, the first $k$ residues are distinct, while the allowed alphabet has size $\tau+1\le k-1$. If $k>p$, the first $p$ residues cover all of $\F_p$, while the alphabet has size $\tau+1\le p-1$. Both cases are impossible.
\end{proof}

\smallskip
\noindent
The following elementary graph lemma makes the construction finite and explicit.

\begin{lemma}[Nearly-regular graphs]\label{lem:nearly-regular-graph}
Let $D\ge0$ and $r\ge D+1$ be integers. There is a simple graph on $r$ vertices with maximum degree at most $D$ and exactly
\begin{equation}\label{eq:graph-edge-count}
\floor{\frac{Dr}{2}}
\end{equation}
edges.
\end{lemma}

\begin{proof}
The case $D=0$ is trivial. Identify the vertices with $\Z/r\Z$. If $D$ is even, join each vertex to the vertices at cyclic distances $1,\ldots,D/2$; this is a $D$-regular graph. If $D$ is odd and $r$ is even, use cyclic distances $1,\ldots,(D-1)/2$ together with the antipodal perfect matching; this is again $D$-regular. These two cases cover $Dr$ even.

\bigskip
\noindent
If $Dr$ is odd, then both $D$ and $r$ are odd, and $r\ge D+2$. Begin with the $(D-1)$-regular circulant graph using cyclic distances $1,\ldots,(D-1)/2$. The edges of cyclic step $(r-1)/2$ form an odd cycle of length $r$ and are disjoint from the existing edge set. Add a maximum matching from that cycle. This adds $(r-1)/2$ edges, leaves one vertex of degree $D-1$, and gives every other vertex degree $D$. The resulting number of edges is
\[
\frac{(D-1)r}{2}+\frac{r-1}{2}=\frac{Dr-1}{2},
\]
as required.
\end{proof}


\begin{theorem}[Graph construction]\label{thm:graph-construction}
Let $p\ge3$ be prime, let $k\ge3$, and put
\begin{equation}\label{eq:digit-capacity}
q=\min\{p,k\}-1,
\qquad
\tau=q-1=\min\{p,k\}-2.
\end{equation}
For every integer $r\ge q-1$, there is a set $A_r$ of distinct positive integers such that
\begin{equation}\label{eq:graph-size}
|A_r|=\floor{\frac{qr}{2}},
\end{equation}
\begin{equation}\label{eq:graph-height}
\max A_r<2p^{r-1},
\end{equation}
and $H(A_r)$ contains no nonconstant $k$-term arithmetic progression.
\end{theorem}

\begin{proof}
Apply Lemma~\ref{lem:nearly-regular-graph} with $D=q-2=\tau-1$ to obtain a simple graph $G$ on vertex set $\{0,1,\ldots,r-1\}$ having maximum degree at most $\tau-1$ and exactly $\floor{(\tau-1)r/2}$ edges. Define
\begin{equation}\label{eq:graph-generators}
A_r=\{p^i:0\le i<r\}
 \cup\{p^i+p^j:\{i,j\}\in E(G)\}.
\end{equation}
All elements in \eqref{eq:graph-generators} are distinct. Moreover,
\[
|A_r|=r+\floor{\frac{(\tau-1)r}{2}}
 =\floor{\frac{(\tau+1)r}{2}}
 =\floor{\frac{qr}{2}}.
\]
In any subset sum of $A_r$, the coefficient of $p^i$ is at most
\[
1+\deg_G(i)\le\tau<p.
\]
Therefore no carry occurs, and every base-$p$ digit lies in $\{0,1,\ldots,\tau\}$. Lemma~\ref{lem:restricted-digits} shows that $H(A_r)$ is $k$-term-progression-free. Finally, every generator is less than $2p^{r-1}$, proving \eqref{eq:graph-height}.
\end{proof}


\end{document}
